\documentclass[UTF8,a4paper,11pt,fontset=fandol]{ctexart}
\usepackage{amsmath,amssymb}
\usepackage[a4paper,margin=2.5cm]{geometry}
\usepackage{hyperref}
\hypersetup{hidelinks}
\setlength{\emergencystretch}{3em}
\title{论文初稿：基于局部区域的晶圆缺陷检索}
\author{}
\date{2026年11月16日}
\begin{document}
\maketitle


本稿对应当前 cc8＋径向位置＋相对面积实现，更新于 2026-11-16。是问题、方法与实验说明的简稿，尚未包含完整相关工作和文献表。本稿与 Markdown 主题页同步。

\section{研究问题}

给定一张含缺陷的晶圆图，从固定候选库中寻找缺陷模式相近的图像。我们关注查询中的缺陷是否被候选覆盖，允许候选有额外缺陷；因此相似度具有方向性。评价同时考察完整标签组合相同的 exact，以及查询标签被覆盖程度的 coverage。后者是晶圆级标签评价，不代表已验证逐区域语义对应。

\section{方法}

在原始 52\(\times\)52 晶圆图的缺陷像素上提取八连通区域，过滤面积小于 5 的区域，最多保留最大的 4 个。以区域中心取固定 52\(\times\)52 窗口，缩放至 96\(\times\)96，分别以所有缺陷、目标区域、有效晶圆构成三通道输入。ResNet18 首层步长设为 1，提取 512 维区域特征。

区域编码器采用自监督 NCE 训练 100 epoch，seed=2026；训练有四种直角旋转，正式推断不使用随机增强或四旋转平均。训练投影维度为 128，检索使用投影前 h512。固定使用训练来源的 \texttt{best\_nce}，增加几何评分不重训或重新选模。

区域对分数同时考虑特征、径向位置和相对面积：

\begin{align*}
s(q,c) &= \operatorname{clip}(\cos(h_q,h_c),0,1)\, e^{-D_{\rm radial}} \\
&\quad\times (0.85+0.15e^{-D_{\rm area}/1.5}), \\
D_{\rm radial} &= \frac{1}{16}\sum_{b=1}^{16}|\operatorname{CDF}_{q,b}-\operatorname{CDF}_{c,b}|, \\
D_{\rm area} &= \frac12 |\log(r_q/r_c)|.
\end{align*}

其中 r 是区域原始缺陷面积占有效晶圆面积的比例。径向项使用归一化半径的 16 格直方图累积分布；中心为有效晶圆包围框中心，半径按有效晶圆最远像素归一化。有效晶圆为原像素值 1 或 2。所有几何量从完整原始 mask 计算，避免裁切和缩放改变位置、尺度。

将区域对分数降序排列，依次选择两端均未占用的配对。最终分数为各查询区域匹配分数的加权平均，权重为原始面积平方根；未匹配区域贡献零。候选额外区域不直接扣分，因此反向查询可能得到不同分数。

\section{实验设置与已得结论}

Mixed-WM38 先精确去重，再按完整标签组合分层，train/valid/test 为 25,193/3,599/7,199。valid 固定 500 个 query，test 使用全部 7,199 个 query；每个 query 检索同集合所有其他图像，Normal 不进入检索库。AP/Recall 分母与 coverage 的理想排序使用完整候选池，报告截断到 K 的指标。

对照包括六种分割、局部/整图融合、占比、角向及匹配策略；外部方法包括 Wu、Yu、WBBS、Nakazawa、WMHD 和 F 对齐的 WaPIRL 整图适配。各方法共享数据与评价协议，保留各自表示，不强行共用区域提取。Yu、Nakazawa 使用监督分类训练，F、WaPIRL 使用自监督训练，其余三种无需训练，比较时注明这一差别。

当前 cc8 在分割对照的 test 前排 exact 指标中较好，但不是所有指标最优；外部比较仍低于 WBBS 和 Nakazawa。融合及角向项改善了检索，最优匹配额外收益较小。完整数值引用 自动结果总览，不在草稿复制易过期的表格。

\section{证据边界}

融合 \(\alpha\)=0.5、角向 \(\gamma\)=0.5 由 valid 的 exact 总体 Precision@10 选择；占比选择零强度。cc8 主方法身份在查看历史 test 后确认，角向＋最优匹配组合也是事后补充，不能写成所有方案都在 test 前确定。自监督训练不等于整个研究过程没有使用标签；valid 选参使用了标签。

目前结果以单个种子为主，尚不能宣称跨种子稳定或差异显著。跨组合评价中的 F 使用旧模型，与当前 cc8 分开报告。方法适配和未公开参数必须按各方法 README 披露，不把当前实现的成绩直接当作原论文成绩。
\end{document}
